\begin{ouverture}[college4]{17}{Le son : fréquence, propagation et écho}{Partie 4 · Des signaux pour observer et communiquer}{Pourquoi voit-on l’éclair avant d’entendre le tonnerre ?}
\objectif{Décrire source vibrante et milieu de propagation.}
\objectif{Distinguer fréquence, hauteur et niveau sonore.}
\objectif{Calculer une distance avec un délai de propagation ou un écho.}
\prerequis{Distance, durée, vitesse et lecture de graphiques.}
\end{ouverture}
\pcPage{Le cours}{Le son : fréquence, propagation et écho}
\phantomsection\label{ess-17}
\essSection{1}{Un son a besoin d’un milieu}
Une source sonore vibre et perturbe un milieu matériel. La perturbation se transmet dans un gaz, un liquide ou un solide ; le son ne se propage pas dans le vide. Les particules du milieu oscillent localement : elles ne voyagent pas toutes de la source à l’oreille. Dans l’air vers $20\,^{\circ}\mathrm C$, on utilise ici $v\approx340\unite{m/s}$. La valeur dépend du milieu et de la température.\par
\pcFigure{figures/c19-signaux.pdf}{Deux signaux périodiques de même amplitude schématique : celui qui présente davantage de motifs par seconde a une fréquence plus élevée.}
\essSection{2}{Fréquence et hauteur}
La fréquence $f$ indique le nombre de motifs périodiques par seconde, en hertz (Hz). Une fréquence plus élevée correspond à un son plus aigu pour les sons simples comparés. L’audition humaine jeune couvre approximativement 20 Hz à 20 000 Hz, avec de fortes variations individuelles. En dessous : infrasons ; au-dessus : ultrasons. Ces mots ne désignent pas des niveaux sonores plus ou moins forts.\par
\clearpage
\essSection{3}{Distance et temps de propagation}
Pour un trajet simple, $d=v\Delta t$. Pour un écho, le son parcourt l’aller et le retour : $2d=v\Delta t$. Le délai entre éclair et tonnerre donne une estimation de distance parce que la lumière arrive bien plus vite que le son. On suppose les deux émissions proches dans le temps et on néglige le délai lumineux pour cette estimation.\par
\pcFigure{figures/c19-echo.pdf}{Un émetteur-récepteur chronomètre le retour ; la distance parcourue par le signal est 2d.}
\essSection{4}{Préserver son audition}
Le niveau sonore, souvent exprimé en décibels (dB), décrit une autre propriété que la fréquence. Le risque auditif dépend notamment du niveau, de la durée et de la répétition des expositions. Une gêne n’est pas un instrument de mesure et l’absence de douleur ne prouve pas l’absence de risque. Réduire le niveau, s’éloigner et limiter la durée sont des moyens de prévention. Aucun essai du livre ne cherche un niveau fort.\par
\begin{pcEncadre}{method}{Un exemple expliqué}Un écho revient $0{,}40\unite s$ après l’émission dans l’air. Le trajet total vaut $340\times0{,}40=136\unite m$. La paroi est à $136/2=68\unite m$. Pour un tonnerre entendu 3 secondes après l’éclair, la distance estimée est au contraire $340\times3=1\,020\unite m$ : on ne divise pas par deux, car c’est un trajet simple.\end{pcEncadre}
\pcPage{Expérimenter et retenir}{Le son : fréquence, propagation et écho}
\phantomsection\label{activite-17}
\begin{pcEncadre}{method}{Investigation en classe}\textbf{Question.} Comment distinguer hauteur et niveau ? Le professeur fournit deux enregistrements modérés ou des tracés, l’un avec davantage de motifs sur une même durée, l’autre avec une amplitude différente. Comparer sans augmenter le volume. Compter les motifs sur une durée donnée et calculer une fréquence. \textbf{Étude de documents :} comparer un délai de tonnerre à un délai d’écho ; dessiner le trajet avant de choisir la formule. \textbf{À rendre :} fréquence, description de la hauteur et distinction avec le niveau.\end{pcEncadre}
\pcSousTitre{Ce que je dois savoir expliquer}\begin{questions}
\item Le son nécessite un milieu matériel ; il ne traverse pas le vide.
\item Une source sonore met un milieu en vibration.
\item La fréquence en Hz caractérise la répétition temporelle.
\item La hauteur et le niveau sonore sont deux propriétés différentes.
\item Un écho correspond à un trajet aller-retour.
\end{questions}
\begin{pcEncadre}{attention}{Une erreur à éviter}Ne pas dire « ultrason = son très fort ». Un ultrason est défini par sa fréquence supérieure à la plage audible humaine usuelle.\end{pcEncadre}
\pcRepere{cor-activite-17}{Pistes pour l’investigation}\pcGrille
\clearpage\markright{Je démarre}
\rubrique{Je démarre}
\par\noindent\begin{minipage}{\linewidth}
\begin{exercice}[1]{Le milieu}
\begin{questions}
\item Le son peut-il se propager dans le vide ?
\item Peut-il traverser de l’eau ?
\end{questions}
\end{exercice}
\end{minipage}
\par\noindent\begin{minipage}{\linewidth}
\begin{exercice}[1]{La source}
\begin{questions}
\item Quelle propriété mécanique permet à une source de produire un son ?
\item Les molécules d’air vont-elles toutes jusqu’à l’oreille ?
\end{questions}
\end{exercice}
\end{minipage}
\par\noindent\begin{minipage}{\linewidth}
\begin{exercice}[1]{L’unité}
\begin{questions}
\item Quelle unité pour la fréquence ?
\item Que signifie une fréquence de 100 Hz pour un signal périodique ?
\end{questions}
\end{exercice}
\end{minipage}
\par\noindent\begin{minipage}{\linewidth}
\begin{exercice}[1]{Trois domaines}
\begin{questions}
\item Classer $10\unite{Hz}$, $1\,000\unite{Hz}$ et $40\unite{kHz}$.
\item Cette classification concerne-t-elle le niveau ?
\end{questions}
\end{exercice}
\end{minipage}
\par\noindent\begin{minipage}{\linewidth}
\begin{exercice}[1]{Le plus aigu}
\begin{questions}
\item Comparer deux sons simples de 200 et 800 Hz. Lequel est le plus aigu ?
\item Est-il nécessairement plus fort ?
\end{questions}
\end{exercice}
\end{minipage}
\clearpage\markright{Je m’entraîne}
\rubrique{Je m’entraîne}
\par\noindent\begin{minipage}{\linewidth}
\begin{exercice}[2]{Une distance simple}
\begin{questions}
\item Dans l’air, un son met $2\unite s$ pour atteindre un auditeur. Calculer la distance avec $340\unite{m/s}$.
\item S’agit-il d’un aller-retour ?
\end{questions}
\end{exercice}
\end{minipage}
\par\noindent\begin{minipage}{\linewidth}
\begin{exercice}[2]{Un écho}
\pcFigureEx{figures/ex-echo-sonore.pdf}{}
\begin{questions}
\item Un écho revient après $0{,}50\unite s$. Calculer la distance à la paroi.
\item Quelle longueur a parcourue le son ?
\end{questions}
\end{exercice}
\end{minipage}
\par\noindent\begin{minipage}{\linewidth}
\begin{exercice}[2]{Compter des motifs}
\pcFigureEx{figures/ex-huit-motifs.pdf}{}
\begin{questions}
\item Un tracé comporte $8$ motifs en $0{,}020\unite s$. Calculer f.
\item Comparer à 200 Hz.
\end{questions}
\end{exercice}
\end{minipage}
\clearpage
\par\noindent\begin{minipage}{\linewidth}
\begin{exercice}[2]{Le tonnerre}
\begin{questions}
\item Un délai de $5\unite s$ sépare éclair et tonnerre. Estimer la distance.
\item Quelle durée a-t-on négligée ?
\end{questions}
\end{exercice}
\end{minipage}
\par\noindent\begin{minipage}{\linewidth}
\begin{exercice}[2]{Dans l’eau}
\pcFigureEx{figures/ex-sonar.pdf}{}
\begin{questions}
\item Un sonar reçoit un écho après $0{,}20\unite s$. On fournit $v=1\,500\unite{m/s}$ dans l’eau. Trouver la profondeur.
\item Pourquoi ne pas utiliser 340 m/s ?
\end{questions}
\end{exercice}
\end{minipage}
\clearpage\markright{Je raisonne}
\rubrique{Je raisonne}
\par\noindent\begin{minipage}{\linewidth}
\begin{exercice}[3]{Un signal plus grand}
\pcFigureEx{figures/ex-meme-periode.pdf}{}
\begin{questions}
\item Deux tracés montrent la même période, mais des amplitudes différentes sur la même échelle. Leur fréquence diffère-t-elle ?
\item Que peut-on dire de leur hauteur si ce sont des sons simples ?
\end{questions}
\end{exercice}
\end{minipage}
\par\noindent\begin{minipage}{\linewidth}
\begin{exercice}[3]{Une information incomplète}
\begin{questions}
\item Un document indique un niveau de 80 dB sans durée ni conditions. Peut-on conclure seul sur tout le risque auditif ?
\item Quels renseignements ajouter ?
\end{questions}
\end{exercice}
\end{minipage}
\par\noindent\begin{minipage}{\linewidth}
\begin{exercice}[3]{Deux trajets confondus}
\begin{questions}
\item Un élève calcule $340\times0{,}30=102\unite m$ comme distance d’une paroi pour un écho. Corriger.
\item Le même produit serait-il pertinent pour un trajet simple de 0,30 s ?
\end{questions}
\end{exercice}
\end{minipage}
\par\noindent\begin{minipage}{\linewidth}
\begin{exercice}[3]{Un délai instrumental}
\begin{questions}
\item Un dispositif ajoute $0{,}010\unite s$ au délai d’écho affiché de $0{,}210\unite s$. Calculer la distance dans l’air après correction.
\item Pourquoi corriger avant de diviser par deux ?
\end{questions}
\end{exercice}
\end{minipage}
\par\noindent\begin{minipage}{\linewidth}
\begin{exercice}[3]{Un contrôle de vraisemblance}
\begin{questions}
\item Un capteur dans l’air annonce un obstacle à 3,4 m avec un écho de 2 s. Est-ce cohérent avec 340 m/s ?
\item Quel délai attendrait-on pour 3,4 m ?
\end{questions}
\end{exercice}
\end{minipage}
\pcPage{Défi du chapitre}{Le son : fréquence, propagation et écho}
\phantomsection\label{defi-ex-17}
\pcDefiTitre{17}{Deux capteurs, deux milieux}
\begin{questions}
\item Un capteur A mesure un écho de $0{,}10\unite s$ dans l’air ; B un écho de même durée dans l’eau. Utiliser respectivement 340 et 1 500 m/s. Calculer les distances.
\item Pourquoi un même délai ne donne-t-il pas la même distance ?
\item Quelle donnée faut-il toujours annoncer dans un compte rendu ?
\end{questions}
\pcGrille
